\section{Aperiodicity and homochirality}\label{sec:aperiodicity}

% ledger: T13, T2, T3, T4, T5, N3; source PROOFS_registered.md R§10; Lean period_halving, all_periods_grid, no_period, sym_card_le_24, r44_einstein (spine); 1711.03401 L8 register
\begin{theorem}[no periods; R\S10; \tier{T1n}, unconditional:
\lean{r44\_einstein}, using \lean{period\_halving} and
\lean{all\_periods\_grid}; terminal arithmetic lemma
\lean{no\_period}: \tier{T1}]\label{thm:noperiod}
Every tiling $T$ by $\solid$ has $\Per(T)=\{0\}$.
\end{theorem}

\begin{proof}
Let $p$ be a translational period of $T$. Registration (\cref{thm:registration}) and the absence
of nonidentity self-isometries (\cref{app:asymmetry}) make $p$ an integer vector: the translation by $p$ carries
a registered tile to a registered tile in the same frame, and the difference of two integer
shifts is an integer vector. The parent partition is translation-covariant (\cref{thm:parent}),
so $D(T+p)=D(T)+p/2$ by \cref{thm:coarsening}; since $T+p=T$, $p/2$ is a period of $D(T)$, which
is again a registered tiling by $\solid$. Iterating, $p/2^n$ is a period of $D^n(T)$ and hence an
integer vector for every $n$, so $p\in\bigcap_n 2^n\Z^3=\{0\}$.
\end{proof}

This is the classical argument that a unique hierarchy forbids periods, in the form the Spectre
paper states \cite{SMKGS24b} and in the form Robinson's tilings first exhibited: a period would
have to be a period of every level, and the levels grow without bound.

\begin{theorem}[finite symmetry; R\S10; \tier{T1n}, unconditional:
\lean{r44\_einstein}; conditional bound
\lean{sym\_card\_le\_24}: \tier{T1}]\label{thm:sym}
Every tiling $T$ by $\solid$ has $|\Sym(T)|\le24$.
\end{theorem}

\begin{proof}
After registration every tile's frame lies in one coset of the $24$-element proper cubic group,
and $\solid$ has no self-isometry, so a symmetry of $T$ is determined by where it sends one
chosen tile: its linear part is one of at most $24$ frames, and, that being fixed, its
translation part is determined as well. The map from $\Sym(T)$ to linear parts has kernel the
translational symmetries, which are trivial by \cref{thm:noperiod}; hence it is injective into a
set of at most $24$ elements.
\end{proof}

\begin{corollary}[strong aperiodicity; a deduction from the two theorems above]\label[corollary]{cor:strong}
No tiling by $\solid$ has a symmetry of infinite order. Hence $\solid$ is a strongly aperiodic
monotile of $\R^3$ in the sense of Mozes and of the hat paper, and \cref{thm:main} holds.
\end{corollary}

\begin{proof}
A finite group has no element of infinite order; existence is \cref{thm:existence}.
\end{proof}

In the taxonomy of Coulbois, Gajardo, Guillon and Lutfalla, \cref{thm:sym} says that $\solid$
is mildly aperiodic. Strong aperiodicity in their sense would require every tiling to have
trivial symmetry group, and $\solid$ does not have that property: the substitution fixed points
of \cref{prop:seedsym} include tilings whose symmetry group has order $8$, as well as tilings
with trivial symmetry group. In particular no screw motion, rational or
irrational, is a symmetry of any tiling by $\solid$, which is the property the biprism lacks.

% ledger: T6, T7, D5; Lean tiling_homochiral, tiling_chirality_corollary (spine); source 2305.17743 L22
\begin{theorem}[homochirality; \lean{tiling\_homochiral}, \lean{tiling\_chirality\_corollary}:
\tier{T1n}, unconditional]\label{thm:homochiral}
In every tiling $T$ by $\solid$ all placements have the same handedness: the determinants of
their linear parts are all $+1$ or all $-1$. Consequently a reflected and an unreflected copy of
$\solid$ never occur in the same tiling, and every tiling is homochiral.
\end{theorem}

\begin{proof}
After the ambient isometry of \cref{thm:registration}, every placement is a proper cubic frame
followed by a translation, so all placements have determinant $+1$ relative to that isometry;
undoing it multiplies all determinants by the same sign. For two tiles $g\solid$ and $h\solid$
the relative motion $hg^{-1}$ then has determinant $+1$, which is the Spectre paper's definition
of a homochiral tiling.
\end{proof}

\begin{corollary}[strictly chiral aperiodic monotile; a deduction at the same verification
boundary]\label[corollary]{cor:chiral}
$\solid$ is a strictly chiral aperiodic monotile of $\R^3$ in the sense of the Spectre paper: it
admits tilings, and every tiling it admits is homochiral and non-periodic.
\end{corollary}

Reflections were never excluded by convention: all $48$ signed frames were admitted as candidate
poses in \cref{sec:companions,sec:registration}, and the $44$ legal contacts turned out proper.
A tiling may be reflected as a whole; what the geometry forbids is a mixture. The
Schmitt--Conway--Danzer biprism, by contrast, needs reflections forbidden for its aperiodicity
\cite{ST12}. The four clauses of \cref{thm:full} are now proved: (1) is \cref{thm:existence},
(2) is \cref{thm:noperiod,thm:sym}, (3) is \cref{thm:homochiral}, and (4) is
\cref{thm:hierarchy}.
